% ============================================================
% 分片：s04b-branes  对应原文 PDF 第 51--56 页（印刷页 51--56）
% 负责章节：4.7 The B-Field、4.8 Tame And Irregular Singularities、4.9 The Multi-Dimensional Case
% 译者：trans-s04   审校状态：初稿
% ============================================================

% ==== 原文 PDF 第 51 页（印刷页 51）====
% （第 51 页上半部分为 4.5--4.6 的结尾，属 s04a-branes 分片；本分片自 4.7 小节标题起。）

\subsection{B-场}

有两处我们遇到了重要的参数，它们不由 A-膜的几何决定，即 $D\subset\widetilde{T}^*Y$ 的极点处的分歧参数 $\alpha^*$，以及式 (4.23) 的扭化类 $\zeta'$。抽象地说，这些参数的任意取值都有意义，都可用来通过求解带奇点的 Hitchin 方程的扭化版本构造扭 D-模。我们想知道在几何 Langlands 的语境中哪些取值是相关的。

对扭化参数，我们可以作出一个基于全纯性的论证，它与我们已就分歧参数所说的非常相似。上同调类 $\chi$，如同扭化参数 $\varrho^*$，必须全纯地依赖于 ${}^L\alpha-i{}^L\gamma$。正如我们刚刚看到的，$\mathrm{Im}\,\chi$ 由 $\mathcal{M}_H$ 中某个 Zariski 开集与余切丛 $T^*\mathcal{M}$ 相差的仿射形变所决定。在 [GW] 中，通过分析 $\mathcal{M}_H$ 的几何，已证明它关于 ${}^L\gamma$ 是线性的，因而 $\mathrm{Re}\,\chi$ 同样关于 ${}^L\alpha$ 是线性的。

% ==== 原文 PDF 第 52 页（印刷页 52）====

然而，我们想更直接地解释 $\zeta'=\mathrm{Re}\,\chi$ 在量子场论中如何起源。到此为止，在我们对辛流形（例如 $X=\widetilde{T}^*Y$）上 A-膜的研究中，我们忽略了 A-模型的一个重要的方面，即 B-场。这样，我们纯粹用 $X$ 的辛形式 $\omega$ 定义了 A-模型。一般地，A-模型全纯地依赖于复化辛形式 $\widehat{\omega}=\omega+iB$。

在没有 B-场时，拉格朗日子流形上的 A-膜由秩 $n$ 向量丛 $V\to L$ 连同平坦酉联络 $A$ 刻画。（为简洁起见，我们在此略去按 $K_L^{1/2}$ 的扭化。）于是，记 $F$ 为 $A$ 的曲率，条件是 $F=0$。然而当 $B\neq0$ 时，该条件被修改为包含一个中心曲率：\footnote{物理学家通常略去右端的因子 $i$，因为他们把 $F$ 定义为取值于厄米而非反厄米 $n\times n$ 矩阵。}
\begin{equation}\tag{4.24}
F=1\otimes iB.
\end{equation}
于是，式 (4.23) 的扭化场 $\zeta'$ 正是 $B$。$B$ 其实只在模去一个正合二形式的加法下才有定义。

% [原文疑误: 「A-branes on a Lagrangian submanifold $X$」中的 $X$ 应为 $L$：紧随其后说「rank $n$ vector bundle $V\to L$」，且 $X$ 已是环境辛流形，若膜支撑在 $X$ 上则拉格朗日条件无意义。]

% [Tutor Note: 物理图像。(4.24) 表明 B-场在 A-模型中的作用是给拉格朗日膜上的平坦酉联络加一个中心（标量）曲率：膜的“丛”不再是平坦的，而是具有曲率 $iB$。由于 $B$ 只模去正合二形式有定义，这一曲率条件在整体上正是 gerbe 联络的语言（见 §10）；在 A-膜的范畴层面，它对应于整个 A-范畴被同一个 gerbe 扭化。]

如 [GW] 所分析，在从 ${}^LG$ 的 B-模型到 $G$ 的 A-模型的对偶下，${}^L\alpha$ 映到 A-模型的一个量子参数 $\eta$，而 $\mathcal{M}_H(G)$ 上的 B-场又关于 $\eta$ 是线性的，其意义如下。$\eta$ 取值于 $\mathfrak{t}$，即 $G$ 的紧实形式的 Lie 代数。$B$ 的上同调类 $[B]$ 是 $H^2(\mathcal{M}_H,\mathbb{R})$ 的元素，后者自然地可等同于 $\mathbb{R}\oplus\mathfrak{t}$。（为简单起见，我们考虑只有一个分歧点的情形；否则有类似的、含更多参数的分析。）相对于此，$[B]=c\oplus\eta$，其中 $c$ 是一个常数，在几何 Langlands 的通常情形为零（而在一个称为量子几何 Langlands 的推广中非零）。

% [Tutor Note: 推导补足。$H^2(\mathcal{M}_H,\mathbb{R})\cong\mathbb{R}\oplus\mathfrak{t}$ 的分解来自 Hitchin 纤维化的结构：常数部分 $\mathbb{R}$ 是整体标量（相当于把 B-场沿整个 $\mathcal{M}_H$ 拉回），$\mathfrak{t}$ 部分则与分歧参数 $\gamma$ 的方向对偶。$c\neq0$ 的情形即“量子”几何 Langlands（在四维规范理论中对应加入 theta 角），此时 Hitchin 纤维化被整体重新标度。]

先考虑 $\mathcal{M}_H^0$ 约化为 $T^*\mathcal{M}$ 的情形。如 [GW] 所示，当 ${}^L\beta={}^L\gamma=0$ 时这确实发生。此时 $Y=\mathcal{M}$ 可作为拉格朗日子流形嵌入。考虑式 (4.23) 中 $\varphi=0$ 的特殊情形。若置 $\zeta=0$ 与 $\zeta'=B$，它就约化为式 (4.24)，因而是描述支撑在 $\mathcal{M}\subset T^*\mathcal{M}$ 上的秩 $n$ A-膜的正确方程。因此 A-膜范畴的扭化参数 $\zeta'$ 正是 $B$，更确切地说正是它的上同调类；于是 $\zeta'$ 与底层 ${}^LG$ 分歧局部系统的分歧参数 ${}^L\alpha$ 一致。$\varphi\neq0$ 的解描述 A-膜从支撑在 $\mathcal{M}$ 上的秩 $n$ 膜形变到支撑在更一般的拉格朗日子流形 $D\subset T^*\mathcal{M}$ 上的秩 1 膜（$D$ 是 $\mathcal{M}$ 的 $n$ 重覆盖）的可能形变。支撑在 $D$ 上的 A-膜以同样的方式被扭化，因为扭化作用于整个 A-范畴。

% [Tutor Note: 易错预警。原文在此处把 (4.23) 中的场写成 $\phi$，把 Higgs 场写成 $\varphi$（“with $\varphi=0$”“with $\varphi\neq0$”），即同一对象在同一段中用了两种字体。按图像照抄；理解为同一 Higgs 场。另请注意：$\zeta'=B$ 的等式是在“$\mathcal{M}_H^0$ 约化为 $T^*\mathcal{M}$”这一特殊情形下论证的，一般情形只保证上同调类相同。]

在几何 Langlands 中，我们真正感兴趣的是 $D\subset T^*\mathcal{M}$ 带极点的情形。然而 $B$ 限制到 $T^*\mathcal{M}$ 时是从 $\mathcal{M}$ 拉回的（模去一个无关的正合二形式），扭化与上一段完全相同。打开 ${}^L\beta$ 与 ${}^L\gamma$（它们是 $\beta$ 与 $\gamma$ 的倍数）会把经典几何从 $T^*\mathcal{M}$ 形变到 $\widetilde{T}^*\mathcal{M}$，但（如 [GW] 所分析）不改变 $B$ 的上同调类，因而并不影响扭化的这一贡献。

% ==== 原文 PDF 第 53 页（印刷页 53）====

总之，在几何 Langlands 的通常情形（带一个分歧点），$\zeta'$ 等于 $0\oplus{}^L\alpha$。带 $k$ 个分歧点时，$H^2(\mathcal{M}_H,\mathbb{R})=\mathbb{R}\oplus(\oplus_{i=1}^{k}\mathfrak{t}^{(i)})$，其中 $\mathfrak{t}^{(i)}$ 是 $\mathfrak{t}$ 的若干副本。每个分歧点有相应的参数 ${}^L\alpha^{(i)}$，类似的论证表明 $\zeta'=0\oplus(\oplus_{i=1}^{k}{}^L\alpha^{(i)})$。

\subsubsection{\texorpdfstring{$\alpha_i^*$}{αi*} 的猜想性描述}

既然已确定 $\zeta'$，我们来讨论参数 $\alpha_i^*$，其中 $i$ 一如既往地标记 $\mathcal{M}$ 中使 Hitchin 纤维「趋于无穷」的除子。

这里关于 B-场的一些一般性说明可能有用。我们所考虑的这一类一般 B-场，在数学上描述为带联络的 $U(1)$-gerbe（类似的对象还将在 §10 中讨论）。这种对象的一个规范变换，是与一个带酉联络 $\nabla$ 的线丛 $\mathcal{L}$ 作张量积。在我们的情形，该 gerbe 其实是平凡的，尽管不是典范地平凡；这就是为什么它的联络可以用上文讨论过的整体定义的二形式 $B$ 来表示。在规范变换下，这个二形式被 $\nabla$ 的曲率平移。

现在设 $\boldsymbol{F}$ 是 Hitchin 纤维化的紧纤维，$\boldsymbol{F}_0$ 是它到 $\widetilde{T}^*\mathcal{M}$ 的限制。$B$ 限制到 $\boldsymbol{F}$ 的上同调类其实为零。这可由 [GW] 中的分析看出，并保证支撑在 $\boldsymbol{F}$ 上的秩 1 A-膜可以存在。这样的 A-膜由给出 gerbe 的一个平坦平凡化 $\mathcal{T}$ 来定义，即相对于该平凡化，联络形式为零。（人们常常假设该 gerbe 有一个先验的平凡化，此时 A-膜由 $\boldsymbol{F}$ 上的平坦线丛定义。然而，内蕴的表述是：A-膜由 gerbe 的平坦平凡化定义。）另一方面，把 $B$ 限制到 $\boldsymbol{F}_0$ 上时，我们可以用一个使 $B$ 为从 $\mathcal{M}$ 拉回的规范。如此 $B$ 完全可能非零，但我们可以取它光滑，从而其范数有界。两个规范之间的差别由带联络 $\nabla$ 的线丛 $\mathcal{L}\to\boldsymbol{F}_0$ 给出。该联络并不平坦，但由于 $B$ 在 $\mathcal{M}$ 上有界，其曲率一致有界。因此 $\mathcal{L}$ 绕 $\boldsymbol{F}\backslash\boldsymbol{F}_0$ 中各除子的单值良定义。我们猜想它们实际上就是分歧参数 $\alpha_i^*$ 的来源。

更确切地说，在 $\mathcal{M}$ 中每个使 Hitchin 纤维 $\boldsymbol{F}$ 趋于无穷的除子 $D_i$ 的邻域内，$\boldsymbol{F}$ 有 $n$ 个分支。因此联络 $\nabla$ 给出一族单值 $\lambda_i^j\in U(1)$，$j=1,\dots,n$，每个分支一个。我们猜想 $\exp(-2\pi\alpha_i^*)$ 属于 $\mathrm{diag}(\lambda_i^1,\dots,\lambda_i^n)$ 的共轭类。这就完成了我们对 $\alpha_i^*$ 的（猜想性）描述。

最后，我们得到与 A-膜 $(\boldsymbol{F},\mathcal{T})$ 相联系的 D-模 $\mathcal{F}$ 的如下描述。由于 $B$ 具有 (1,1) 型，$\mathcal{L}\to\boldsymbol{F}_0$ 的曲率是 (1,1) 型的，故联络 $\nabla$ 赋予 $\mathcal{L}$ 一个全纯结构。我们把 $\mathcal{L}$ 延拓为 $\boldsymbol{F}$ 上的全纯线丛，其方式是：称 $\mathcal{L}\to\boldsymbol{F}_0$ 的全纯截面 $\psi$ 在 $\boldsymbol{F}$ 上全纯，若其模长 $|\psi|$（用 $\mathcal{L}$ 的酉结构计算）有界。此外，我们称 $\psi$ 在 $\boldsymbol{F}\backslash\boldsymbol{F}_0$ 的不可约分支 $D'$ 附近以 $\alpha$ 阶消没，若在 $D'$ 附近有 $|\psi|\sim|z|^{\alpha}$，其中 $z$ 是 $D'$ 附近的法坐标。截面的消没阶的这一刻画是良定义的，因为 $\mathcal{L}$ 的曲率有界。它赋予 $\mathcal{L}\to\boldsymbol{F}$ 沿 $\boldsymbol{F}\backslash\boldsymbol{F}_0$ 的抛物结构。

% ==== 原文 PDF 第 54 页（印刷页 54）====
于是 $E=\pi_*(\mathcal{L})$（其中 $\pi:\boldsymbol{F}_0\to\mathcal{M}$ 是投影）在除子 $D_i$ 处有自然的抛物结构。它还以通常的方式带有自然的 Higgs 场 $\varphi$。

我们的提议是：$\mathcal{F}$ 限制到除子 $D_i$ 的补集上，是由抛物 Higgs 丛 $(E,\varphi)$ 作非阿贝尔 Hodge 变换得到的局部系统。

% [Tutor Note: 推导补足。这里的构造是“抛物非阿贝尔 Hodge 对应”的标准用法：先用 $B$ 的 (1,1) 型给 $\mathcal{L}$ 一个全纯结构，再用 $|\psi|$ 有界把 $\mathcal{L}$ 延拓到 $\boldsymbol{F}$（即取 $\boldsymbol{F}_0$ 上的“有界截面”），用 $|\psi|\sim|z|^{\alpha}$ 定义沿边界除子的消没阶，从而得到沿 $\boldsymbol{F}\backslash\boldsymbol{F}_0$ 的抛物结构；对 $(E,\varphi)$ 作非阿贝尔 Hodge 变换（即解 Hitchin 方程）便得到 $\boldsymbol{F}$ 的补集上的局部系统。曲率有界保证 $|\psi|$ 的增长阶与法坐标选取无关，这是抛物结构良定义的关键。]

\subsection{驯顺奇点与不规则奇点}

迄今我们只考虑了刻画拉格朗日曲线 $D\subset T^*Y$ 奇性行为的指数 $d_i$
\begin{equation}\tag{4.25}
s_i(u)\sim c_i(u-u_0)^{-d_i},\qquad i=1,\dots,n
\end{equation}
等于 1 的情形。这里我们将简要地放宽这一假设。

一种情形是指数为正（于是 $D$ 确实趋于无穷）但小于 1。这只可能在 $D$ 通过添加一个或多个无穷远点而完备化后于 $u_0$ 处分歧时发生。例如取 $u_0=0$，考虑谱曲线 $\det(s-\varphi)=0$，其中 $\det\varphi=1/u$。此时谱曲线的方程是 $s^2+u^{-1}=0$，分支为 $s=\pm(-u)^{1/2}$，故指数 $d_i$ 等于 $1/2$。$\varphi$ 的一个合适选择是
\begin{equation}\tag{4.26}
\varphi=du\begin{pmatrix}0&u^{-1}\\-1&0\end{pmatrix}.
\end{equation}
$\varphi$ 有极点，但极点的留数是幂零的，故 $\det\varphi$ 只有一个单极点，而不是二阶极点。此时适用 [Sim1] 的分析；它对应于 $\beta^*=\gamma^*=0$，如 [GW] §3.3 所回顾的。由于 $\gamma^*=0$，由支撑在 $D$ 上的 A-膜在 $Y$ 上产生的局部系统的单值其模为 1。

相反的情形是指数 $d_i$ 大于 1。此时 $\varphi$ 在 $u_0$ 处有高于一阶的极点。对于 $\varphi$ 有这种不规则极点的 Higgs 丛，其 Hitchin 方程已在 [BB] 中分析。相应的局部系统在 $u_0$ 处有不规则奇点。联络的不规则部分可以直接由 Higgs 场的奇点读出，这与正则情形中参数 $\alpha,\beta,\gamma$ 之间微妙的相互配合形成对照。非正式的说明见 [Wi2]。

% [Tutor Note: 易错预警。$0<d_i<1$ 并不表示奇性更弱，而是表示 $D$ 的完备化在 $u_0$ 处分歧（$D$ 仍趋于无穷）；例 (4.26) 展示的是：$\varphi$ 有极点但留数幂零，于是 $\det\varphi$ 的极点阶数降低——这是“驯顺而分歧”的典型机制，也正是 [Sim1] 的框架所能处理的。$d_i>1$ 才是真正的不规则奇点，此时单值的模不再为 1，且不规则部分不经过 $\alpha,\beta,\gamma$ 的微妙配合，可直接读出。]

\subsection{多维情形}

对几何 Langlands，我们需要把上述一切推广到 $\widetilde{T}^*Y$ 中的 A-膜，其中 $Y$ 的维数大于 1。Higgs 丛与局部系统之间的对应可以推广到更高维 [Sim2, Sim3]，而这也推广到带极点的 Higgs 丛 [Bi, Mo]。允许极点这一推广对几何 Langlands 至关重要，因为 Hitchin 纤维化的纤维 $\boldsymbol{F}$ 限制到 $T^*\mathcal{M}$ 时本质上总有极点。不过这里我们只考虑无极点的情形，这已足以说明基本想法。

复流形 $Y$（维数为 $n$）上的 Higgs 丛定义为对 $(E,\varphi)$，其中 $E\to Y$ 是全纯向量丛，而 $\varphi:E\to E\otimes\Omega^1(Y)$ 是全纯映射，满足 $\varphi^2=0$（$\varphi^2$ 由复合 $E$ 的自同态并利用杯积 $\wedge^2\Omega^1\to\Omega^2$ 定义）。最后一个条件在曲线情形是平凡的。基本结果是：若 $Y$ 是 Kähler 的，则 $Y$ 上 $E$ 的有理 Chern 类为零的 Higgs 丛与 $Y$ 上的局部系统之间存在自然的对应。该对应通过求解 Hitchin 方程的多维推广来实现。（若 $Y$ 是射影的，这一操作等价于在 $Y$ 中的每条曲线上同时求解 Hitchin 方程。）

% ==== 原文 PDF 第 55 页（印刷页 55）====
% [跨页接续: 原书第 54/55 页，「a natural correspondence between Higgs bundles on $Y$ ...」跨页。]
Higgs 丛 $(E,\varphi)$ 与装备了线丛 $\mathcal{L}$ 的谱簇 $D\subset T^*Y$ 之间也有自然的对应。记 $\pi:T^*Y\to Y$ 为投影，该对应在一个方向上由 $E=\pi_*(\mathcal{L})$ 与 $\varphi=\pi^*(z)$ 给出，其中 $z$ 在余切丛的纤维上是线性的。

我们希望把这一对应陈述为关于支撑在光滑嵌入的紧拉格朗日子流形 $D\subset T^*Y$ 上的秩 1 A-膜的结果。（$\dim Y>1$ 时这种情形的实例其实极为稀少，这也是必须推广到允许 $D$ 有极点的另一个理由。）这样的秩 1 A-膜存在，当且仅当 $D$ 的典范丛具有平方根（这与注 4.1 有关）。为指定这样的膜，我们在 $D$ 上选取一个平坦 $\mathrm{Spin}_c$ 结构，我们并不十分内蕴地把它描述为线丛 $K_D^{1/2}\otimes\mathcal{N}$ 的选择，其中 $\mathcal{N}$ 的第一 Chern 类为零。然后我们定义\footnote{若 $K_Y$ 没有全局平方根，则 $E$ 必须理解为扭向量丛，被某个 gerbe 扭化。无论如何，那是最内蕴的表述。}向量丛 $E\to Y$，使 $K_Y^{1/2}\otimes E=\pi_*(K_D^{1/2}\otimes\mathcal{N})$，换言之 $E=\pi_*(\mathcal{L})$，其中 $\mathcal{L}=K_D^{1/2}\otimes\pi^*(K_Y^{-1/2})\otimes\mathcal{N}$。这延拓为 Higgs 丛 $(E,\varphi)$，其中 $\varphi=\pi^*(z)$。$E$ 的秩与覆盖 $\pi:D\to Y$ 的次数 $n$ 一致。

Higgs 丛 $(E,\varphi)$ 将对应于 $Y$ 上的局部系统，若 $E$ 的 Chern 特征（我们记作 $\mathrm{ch}(E)$）等于 $n$。一个通过 Riemann--Roch 的论证表明这总是成立的。陈述这一论证时，我们用到 $\widehat{A}$ 类，它与复流形 $X$ 的 Todd 类的关系是 $\widehat{A}(X)=\mathrm{Td}(X)\mathrm{ch}(K_X^{1/2})$。$\widehat{A}$ 类对任何实向量丛 $V$ 都有定义，并具有性质 $\widehat{A}(V)=\widehat{A}(V^*)$，其中 $V^*$ 是 $V$ 的对偶。

% [Tutor Note: 易错预警。记号是 $\widehat{A}$（A-hat，读作“A 帽”），不是 $A$；文字层常把它渲染成“bA”或“A”，务必按图像核对。$\widehat{A}(X)=\mathrm{Td}(X)\mathrm{ch}(K_X^{1/2})$ 中 $K_X^{1/2}$ 的存在性是形式上的（取形式平方根），故 $\widehat{A}$ 对任何复流形都有定义，而 $\mathrm{Td}$ 只用到切丛。]

Riemann--Roch 公式在此情形说的是
\begin{equation}\tag{4.27}
\pi_*\bigl(\mathrm{Td}(D)\mathrm{ch}(K_D^{1/2})\bigr)=\mathrm{Td}(Y)\mathrm{ch}(K_Y^{1/2})\mathrm{ch}(E).
\end{equation}
我们用到了 $c_1(\mathcal{N})=0$ 故 $\mathrm{ch}(\mathcal{N})=1$ 这一事实。用 $\widehat{A}$ 类来说，这意味着
\begin{equation}\tag{4.28}
\pi_*(\widehat{A}(D))=\widehat{A}(Y)\mathrm{ch}(E).
\end{equation}
因此 $\mathrm{ch}(E)$ 等于 $n$，当且仅当 $\pi_*(\widehat{A}(D))=n\widehat{A}(Y)$。由于这是有理上同调中的陈述，它等价于证明
\begin{equation}\tag{4.29}
2\pi_*(\widehat{A}(D))=2n\widehat{A}(Y).
\end{equation}
现在按定义 $\widehat{A}(D)$ 是 $D$ 的切丛 $TD$ 的 $\widehat{A}$ 类。故 $2\widehat{A}(D)=\widehat{A}(TD\oplus TD)$。由于 $D\subset T^*Y$ 是拉格朗日的，$D$ 的法丛与切丛同构（作为实向量丛）。它们的直和是 $T(T^*Y)|_D$，即 $T^*Y$ 的切丛 $T(T^*Y)$ 到 $D$ 的限制。于是我们可以把 $TD\oplus TD$ 换成 $T(T^*Y)|_D$，并把 $2\pi_*(\widehat{A}(D))$ 解释为 $\pi_*(\widehat{A}(T(T^*Y)|_D))$。由于 $Y\subset T^*Y$ 也是
% ==== 原文 PDF 第 56 页（印刷页 56）====
% [跨页接续: 原书第 55/56 页，「As $Y\subset T^*Y$ is also Lagrangian ...」跨页。]
拉格朗日的，我们同样可以把 $2\widehat{A}(Y)$ 解释为 $\widehat{A}(T(T^*Y)|_Y)$。于是我们想要的公式是
\begin{equation}\tag{4.30}
\pi_*\bigl(\widehat{A}(T(T^*Y)|_D)\bigr)=n\widehat{A}(T(T^*Y)|_Y).
\end{equation}

由于 $T^*Y$ 可缩到 $Y$，任何向量丛 $V\to T^*Y$ 都同构于其到 $Y$ 的限制从 $Y$ 拉回，即 $V\cong\pi^*(V|_Y)$。取 $V=T(T^*Y)$，得 $T(T^*Y)=\pi^*(T(T^*Y)|_Y)$。于是特别地 $T(T^*Y)|_D$ 是 $\pi^*(T(T^*Y)|_Y)$ 到 $D$ 的限制。因此式 (4.30) 的左端是 $\pi_*(\widehat{A}(\pi^*(T(T^*Y)|_Y)))$。一般地，对任何映射 $\pi:X\to Y$ 与实向量丛 $V\to Y$，有 $\widehat{A}(\pi^*(V))=\pi^*(\widehat{A}(V))$。故 (4.30) 的左端是 $\pi_*(\pi^*(\widehat{A}(T(T^*Y)|_Y)))$。对 $n$ 次映射 $\pi:D\to Y$，复合 $\pi_*\pi^*$ 在上同调上的作用是乘以 $n$，于是式 (4.30) 的正确性便清楚了。

% [Tutor Note: 推导补足。整段论证的骨架是：把 $2\widehat{A}$ 写成 $\widehat{A}$ 的“二倍丛”形式（$2\widehat{A}(D)=\widehat{A}(TD\oplus TD)$），再用 $D$ 与 $Y$ 都是拉格朗日这一事实把 $TD\oplus TD$ 与 $T(T^*Y)|_D$、$TY\oplus TY$ 与 $T(T^*Y)|_Y$ 对应起来，最后用 $T^*Y$ 可缩到 $Y$（故任何丛都是拉回）与 $\pi_*\pi^*=n$（对 $n$ 次覆盖）收尾。注意 $\widehat{A}$ 对拉回自然、对偶丛不变（$\widehat{A}(V)=\widehat{A}(V^*)$）这两条性质是全部计算的支点。]

与 Higgs 丛 $(E,\varphi)$ 相对应的 $Y$ 上的局部系统，就是与我们的出发点——支撑在拉格朗日子流形 $D\subset T^*Y$ 上的秩 1 A-膜——相对应的所需 D-模。这是一个非常特殊的情形，因为我们假设了 $D$ 紧。在大多数与几何 Langlands 纲领相关的例子中（例如本节前面考虑过的那些），并非如此，相应的 D-模并不由整个 $Y$ 上的局部系统代表，而只由 $Y$ 的某个开子集上的局部系统代表。在 $\dim Y=1$ 的情形，我们在 §4.5--4.7 中已解释了如何在 $Y$ 的开子集上构造相应的局部系统。类似的情形有望在多维情形成立。关于所需的部分分析，见 [Bi], [Mo]。

% [Tutor Note: 物理图像。$D$ 紧的秩 1 A-膜给出定义在整个 $Y$ 上的局部系统，这是极为特殊的情形；几何 Langlands 中真正出现的 Hitchin 纤维（如 §3、§4.5 的例子）不紧，只能给出 $Y$ 的某个开集上的局部系统，其单值在边界除子处有极点，这正是 §4.4--4.7 的整套扭化/极点机制所要处理的内容。]
